\section{Introduction and falsifiable memory question}

The problem of distinguishing fractional, delayed, and latent ecological memory is well-posed only when complexity restrictions (finite latent dimension) are imposed; otherwise, the finite-horizon discrimination problem is ill-posed due to the no-free-lunch theorem.

\textbf{Evidence:} Theorem~T10 (no-free-lunch result) and the minimal viable product (MVP) setup with fixed maximum latent dimension \(m\) (see~\textsc{Idea}~\S~4, \texttt{source\_pack/README.md}).

\textbf{Claim type:} Inference (based on theorem T10).

\textbf{Claim state:} DRAFTING (requires no simulation, permitted after analytical work alone).

This paper establishes a rigorous framework for optimal experimental design to distinguish three mechanistic models of ecological memory: fractional (Caputo derivative of order \(\alpha\)), delayed (discrete time-delay), and latent (finite-dimensional linear latent state of fixed maximal dimension \(m\)).  The analysis proceeds around a certified strong-Allee predator--prey equilibrium, for which we obtain exact linearized formulas (Theorems~T1 and T2), prove structural non-equivalence of the fractional model with finite-dimensional latent and delay alternatives (Theorems~T4 and T5), establish frequency-domain signatures that separate the models (Theorems~T6 and T7), construct finite-horizon exponential approximations with explicit error bounds (Theorems~T9 and T10), derive optimal input and sampling designs under energy and safety constraints (Theorems~T11--T13), and establish the existence of safe, informative perturbations (Theorems~T16 and T17).  A simulation benchmark validates the analytical results and compares excitation waveforms under a common energy budget (Section~\ref{sec:simulation}).  The paper concludes with a implementable laboratory protocol or calibrated digital twin for applying the design to real ecological systems (Section~\ref{sec:protocol}).
\end{document}